% Generated draft: requires mathematical and source review.
\documentclass[11pt]{article}
\usepackage[margin=25mm]{geometry}
\usepackage{fontspec}
\setmainfont{Latin Modern Roman}
\usepackage{amsmath,amssymb,mathtools,unicode-math}
\setmathfont{Latin Modern Math}
\usepackage{enumitem,needspace,xurl,hyperref}
\setlength{\parindent}{0pt}
\setlength{\parskip}{5pt}
\setlength{\emergencystretch}{3em}
\newcommand{\N}{\mathbb{N}}
\newcommand{\Z}{\mathbb{Z}}
\newcommand{\Q}{\mathbb{Q}}
\newcommand{\R}{\mathbb{R}}
\newcommand{\C}{\mathbb{C}}
\newcommand{\F}{\mathbb{F}}
\newcommand{\A}{\mathbb{A}}
\newcommand{\PP}{\mathbb{P}}
\newcommand{\E}{\mathbb{E}}
\newcommand{\eps}{\varepsilon}
\newcommand{\dd}{\,\mathrm{d}}
\begin{document}
\section*{UnsolvedMath ID: 801151}
\textbf{Source problem code:} catalogue entry
\subsection{Definitions and mathematical statement}
The supplied problem asks: given a Nijenhuis operator \(L\), describe all of its local conservation laws. The excerpt does not define ``conservation law'' or specify the manifold category and regularity, so a fully determined classification problem requires these qualifications.

For a conditional smooth real formulation, let \(M\) be a smooth real manifold and let \(L\colon TM\to TM\) be a smooth bundle endomorphism. Its Nijenhuis torsion is
\[
N_L(X,Y)=[LX,LY]-L[LX,Y]-L[X,LY]+L^2[X,Y],
\]
for smooth vector fields \(X,Y\). The endomorphism \(L\) is a Nijenhuis operator if \(N_L(X,Y)=0\) for all such \(X,Y\). Its action on a one-form \(\alpha\) is defined by
\[
(L^*\alpha)(X)=\alpha(LX).
\]
If a scalar conservation law is intended to mean a function \(f\) satisfying \(\mathrm d(L^*\mathrm df)=0\), then the requested problem becomes: for every point \(p\in M\), determine all local solutions, on sufficiently small open neighborhoods \(U\) of \(p\), of
\[
\mathcal C_L(U)=\bigl\{f\in C^\infty(U,\R):\mathrm d(L^*\mathrm df)=0\bigr\}.
\]
If conservation laws are instead represented by one-forms, the corresponding possible target is
\[
\bigl\{\alpha\in\Omega^1(U):\mathrm d\alpha=0,\quad \mathrm d(L^*\alpha)=0\bigr\}.
\]
These are explicitly conditional interpretations: the supplied excerpt does not select either definition, specify an equivalence relation, or state whether constants are to be identified. It also does not specify boundary conditions or a smooth, analytic, or holomorphic category.

The source additionally asks the analogous question for arbitrary operators, without the Nijenhuis assumption, and refers to a discussion of superintegrable potentials that is not included in the supplied record.
\subsection{Short English statement}
Describe every local conservation law associated with a given Nijenhuis operator \(L\). The excerpt leaves the meaning of conservation law and the geometric category unspecified; the displayed differential equations are conditional formulations requiring review. The source also mentions the analogous question for operators that are not Nijenhuis.
\subsection{Sources}
\begin{itemize}
\item Local UnsolvedMath registry, permanent ID \(801151\), source problem code \texttt{catalogue entry}, title ``Nijenhuis Geometry and Integrable Systems --- Problem 4.'' The registry attributes the problem to Bolsinov, Konyaev, and Matveev, with proposed year 2024, and supplies the status label ``open''; this label is not an independently verified literature assessment.
\item Bolsinov, Konyaev, Matveev, \emph{Research Problems on Nijenhuis Geometry and Integrable Systems} (2024), Problem 4. Supplied source URL: \url{https://arxiv.org/abs/2410.04276}. The record reports access on 2026-10-16 and extraction from arXiv LaTeX; these are supplied provenance details, not a claim of independent source access.
\end{itemize}
\end{document}
